% =============================================================================
% chapter4.tex — System Architecture and Design Decisions
% Chapter 4 of the Codara PFE Report
% =============================================================================

\chapter{System Architecture and Design Decisions}
\label{ch:architecture}

% ---------------------------------------------------------------------------
\section{Introduction}
\label{sec:ch4-intro}
% ---------------------------------------------------------------------------

Architecture is not chosen in a vacuum: every decision in this chapter was
preceded by a concrete experiment, a benchmark run, or a measured failure. The
philosophy that guided the design of \codara{} is best expressed as
\emph{empirical selection over theoretical preference}: when two frameworks
appeared equally suitable on paper, we built a prototype with each, measured
the result on a representative SAS workload, and chose the one that performed
better. This section documents that process with the transparency needed for
scientific reproducibility.

The chapter proceeds as follows. Section~\ref{sec:tech-selection} describes
the technology selection process for the four most consequential architectural
choices: orchestration framework, LLM provider and routing strategy, embedding
model and vector store, and formal verification tool. Section~\ref{sec:pipeline-design}
presents the eight-node LangGraph pipeline in detail, explaining the composite
node design pattern and the responsibilities of each node group.
Section~\ref{sec:data-architecture} documents the four-store data architecture
(SQLite, LanceDB, DuckDB, Redis) and the Azure cloud storage layer.
Section~\ref{sec:api-design} covers the API and service-layer design.
Section~\ref{sec:corpus-construction} describes the construction of the
gold-standard evaluation corpus.

% ---------------------------------------------------------------------------
\section{Technology Selection Through Empirical Evaluation}
\label{sec:tech-selection}
% ---------------------------------------------------------------------------

\subsection{Multi-Agent Orchestration Framework Selection}
\label{subsec:framework-selection}

\textbf{The technology.} Three multi-agent orchestration frameworks were
evaluated: \textbf{LangGraph}~[1], \textbf{CrewAI}~[2], and
\textbf{AutoGen}~[3].

LangGraph models workflows as directed graphs over a shared
\texttt{TypedDict} state; each node is a Python coroutine. It is designed for
\emph{deterministic} pipelines where the execution order is known at design
time and the state transitions are explicit.

CrewAI takes an agent-role metaphor: you define ``agents'' with natural-language
role descriptions and ``tasks'' that the agents execute in sequence or in
parallel. It is designed for collaborative, loosely-coupled workflows where
agents communicate through natural language.

AutoGen focuses on conversational multi-agent patterns: agents exchange
messages, including code, in a shared thread, and a designated ``orchestrator''
agent can spawn sub-conversations. It is best suited for open-ended code
generation and debugging scenarios where the workflow is not known in advance.

\textbf{Our evaluation.} We built a three-node prototype (parse, translate,
validate) with each framework and ran it against a 10-file SAS corpus.

[TABLE 4.1 --- Caption: ``Orchestration framework evaluation: LangGraph vs
CrewAI vs AutoGen on the three-node prototype.'' --- Columns: Framework,
State Typing, Checkpoint Support, Async Native, Deterministic Topology,
Lines of Code (3-node), Observed Latency (10-file corpus).]

\textbf{Results and choice.} CrewAI's agent-role abstraction introduced
non-determinism that made debugging difficult: the same SAS file could take
different translation paths on different runs depending on how the LLM
interpreted the role prompts. AutoGen's conversational model was powerful for
open-ended code generation but added latency (each agent ``turn'' requires
an LLM call) that was unnecessary for a pipeline with a fixed execution graph.

LangGraph emerged as the clear choice for \codara{}. Its explicit
\texttt{StateGraph} compilation enforces the pipeline topology at startup,
catching integration errors before any LLM call is made. Its native Redis
checkpoint interface enabled fault recovery with minimal additional code.
And its \texttt{async def} node model allowed concurrent LLM calls within the
translation node via \texttt{asyncio.gather}, reducing translation latency
by approximately 40\% compared to sequential calls. All eight pipeline nodes
were implemented as LangGraph coroutines.

\subsection{LLM Benchmarking and Provider Selection}
\label{subsec:llm-benchmarking}

\textbf{The technology.} LLM selection is arguably the most consequential
single decision in a system of this kind: the quality of the translation is
bounded by the capability of the underlying model. We evaluated three provider
tiers: \textbf{Azure OpenAI} (GPT-4o and GPT-4o-mini), \textbf{Ollama}
(local/cloud-proxied open-weight models), and \textbf{Groq} (LLaMA-3.3-70b
via the Groq fast-inference API).

\textbf{The benchmark.} We designed a ``torture test'' SAS fixture
(\texttt{backend/tests/fixtures/torture\_test.sas}) comprising ten blocks
that represent the hardest migration patterns encountered in practice:

\begin{enumerate}
  \item RETAIN + BY-group FIRST./LAST. accumulation
  \item Missing-value arithmetic (SAS \texttt{.} propagation)
  \item PROC SQL correlated subquery
  \item Macro with \texttt{\%DO} loop
  \item PROC MEANS CLASS OUTPUT statement
  \item PROC SORT NODUPKEY deduplication
  \item Hash object key-value lookup
  \item Multi-level nested macro expansion
  \item PROC TRANSPOSE column reshaping
  \item Complex WHERE with FORMAT inline conversion
\end{enumerate}

Each model was asked to translate all ten blocks, and the translations were
evaluated on: syntax validity (Python \texttt{compile()}), execution without
error on a synthetic DataFrame, and semantic correctness (SemantiCheck Score).

[TABLE 4.2 --- Caption: ``LLM torture-test benchmark results across three
provider tiers. Score out of 10.'' --- Columns: Model, Provider, Syntax
Valid (/10), Execution OK (/10), SemantiCheck VERIFIED or LIKELY\_CORRECT (/10),
Total Score (/10), Avg Latency (s/block).]

\textbf{Evolution of the LLM chain.} The provider selection underwent three
distinct phases, each driven by concrete benchmark evidence:

\begin{itemize}
  \item \textbf{Weeks 1--8}: Groq LLaMA-3.3-70b was the primary translation
    model; Azure GPT-4o was an aspirational fallback that required API access
    not yet provisioned. Ollama ran a local 8B model for LOW-risk partitions.

  \item \textbf{Week 9 (Azure migration)}: Once the Azure OpenAI deployment
    was operational, the chain was restructured to: Azure GPT-4o (primary,
    MOD/HIGH) $\rightarrow$ Azure GPT-4o-mini (LOW) $\rightarrow$ Groq LLaMA
    (cross-verifier). Azure became the primary model because its GPT-4o
    outperformed LLaMA-3.3-70b by 8 percentage points on translation success
    rate. This state persisted through Weeks 9--14 (v3.0).

  \item \textbf{April 6, 2026}: Ollama \texttt{minimax-m2.7:cloud} was
    benchmarked on the torture test and scored \textbf{10/10} in 213 seconds
    --- outperforming GPT-4o on 3 of the 10 blocks. This was a significant
    result because it was an open-weight cloud-proxied model, potentially
    free to use. Following this result, \texttt{nemotron-3-super:cloud} was
    adopted as the Ollama production model (nemotron had better concurrency
    characteristics for pipeline workloads). The chain order was accidentally
    transposed during this update, placing Nemotron as the primary provider.

  \item \textbf{April 20, 2026 (fix)}: The chain order bug was identified and
    corrected. The final production chain is:
    \textbf{Azure GPT-4o} (primary, all risk levels) $\rightarrow$
    \textbf{Ollama \texttt{nemotron-3-super:cloud}} (warm fallback) $\rightarrow$
    \textbf{Groq LLaMA-3.3-70b} (last-resort + independent cross-verifier).
    Azure is kept as primary for production reliability and determinism;
    Ollama is warm fallback rather than experimental.
\end{itemize}

[FIGURE 4.1 --- Caption: ``LLM fallback chain in \codara{} v3.1.0: primary
Azure GPT-4o with Ollama and Groq fallbacks, circuit breakers, and rate
limiters.'' --- Suggested: flow diagram showing the three LLM boxes with
circuit breaker symbols, rate limiter annotations, and the PARTIAL outcome
when all fail.]

\textbf{Circuit breakers and rate limiters.} Two resilience patterns protect
the LLM call chain. A circuit breaker~[4] monitors consecutive failures per
provider and ``opens'' (stops sending requests) after a threshold is reached,
preventing cascading failures when a provider is degraded: Azure opens after
5 failures (60-second cool-down); Groq opens after 3 failures (120-second
cool-down). Rate limiters cap concurrent requests: 10 for Azure, 3 for Groq
(matching Groq's free-tier concurrency limit). Both are implemented as async
semaphores in \texttt{partition/utils/retry.py}.

\subsection{Embedding Model and Vector Store Selection}
\label{subsec:embedding-selection}

\textbf{The technology.} Three embedding models were considered for the
knowledge-base retrieval component: \textbf{Nomic Embed v1.5}~[5] (768-dim,
Apache-licensed, optimised for code-mixed text), \textbf{OpenAI
text-embedding-3-small} (1536-dim, API-based), and \textbf{all-MiniLM-L6-v2}
(384-dim, fast but lower quality).

For the vector store, we evaluated \textbf{LanceDB}~[6] (columnar, embedded,
no server required), \textbf{Chroma}~[7] (server-based, Python-native), and
\textbf{Weaviate}~[8] (cloud-native, production-oriented).

\textbf{Our evaluation.} The evaluation measured retrieval quality (hit-rate@5
on a 50-query evaluation set drawn from the gold-standard corpus), cold-start
latency (time from import to first query), and operational simplicity
(does it require a separate server process?).

[TABLE 4.3 --- Caption: ``Embedding model and vector store evaluation.'' ---
Columns: Model/Store, Dimensions, Hit-Rate@5, Cold-Start (s), Requires Server,
Notes.]

\textbf{Results and choice.} Nomic Embed v1.5 + LanceDB emerged as the
preferred combination. Nomic's 768-dimensional embeddings produced the highest
hit-rate@5 (0.84 vs 0.79 for MiniLM) while remaining fully local, eliminating
API cost and latency for every retrieval call. LanceDB's embedded architecture
--- no server process, files on local disk --- matched the project's
operational simplicity requirement. Its IVF indexing~[9] provides
sub-millisecond approximate nearest-neighbour search over the full 330-entry
knowledge base, and its columnar format supports efficient filtering by
metadata fields (complexity tier, failure mode, verification status) without
loading all embeddings into memory.

The 16-field LanceDB schema includes: \texttt{sas\_snippet} (source text),
\texttt{python\_translation} (target text), \texttt{embedding} (768-dim float
array), \texttt{complexity\_tier}, \texttt{failure\_mode}, \texttt{verified}
(boolean), \texttt{verification\_score}, \texttt{version}, \texttt{superseded\_by},
and seven auxiliary metadata fields.

\subsection{Formal Verification Tool Selection}
\label{subsec:z3-selection}

\textbf{The technology.} Three verification approaches were evaluated: the
\textbf{Z3 SMT solver}~[10] (Microsoft Research, open source), Python-level
\textbf{property-based testing} with Hypothesis~[11], and manual
\textbf{differential testing} against a fixed set of reference inputs.

Z3 is an industrial-strength SMT solver that decides the satisfiability of
first-order logic formulae over arithmetic, bitvector, and array theories.
It provides formal proofs (UNSAT = equivalent for all inputs) rather than
heuristic evidence.

Hypothesis generates random inputs that satisfy property specifications,
providing statistical confidence but not formal guarantees.

Differential testing runs both the SAS-semantics reference implementation and
the generated Python on the same inputs and compares outputs; it is simple to
implement but provides no coverage guarantee.

\textbf{Results and choice.} Z3 was selected as the primary verifier for its
unique property: when it reports UNSAT, the equivalence is \emph{proved}, not
merely likely. For the 30\% of SAS constructs that encode cleanly as linear
arithmetic or Boolean formulae (RETAIN accumulations, WHERE filters, PROC SORT
NODUPKEY, assignment chains), Z3 provides the strongest possible correctness
guarantee.

Hypothesis was retained as a complementary approach for constructs outside Z3's
encoding scope: its randomised inputs form the basis of the CDAIS framework
described in Section~\ref{sec:cdais-impl}. Differential testing is used in the
regression runner (Section~\ref{sec:regression-impl}).

% ---------------------------------------------------------------------------
\section{The 8-Node LangGraph Pipeline}
\label{sec:pipeline-design}
% ---------------------------------------------------------------------------

\subsection{State Model and Composite Node Design Pattern}
\label{subsec:state-model}

The LangGraph pipeline operates over a shared \texttt{PipelineState}
\texttt{TypedDict} that is passed through all eight nodes. The state schema
includes: the list of input files and their metadata, the list of
\texttt{PartitionIR} objects (the core unit of work, each identified by a UUID
\texttt{block\_id}), the RAPTOR tree structure, the risk-level assignments,
the translated Python strings per partition, the verification results, and
accumulated warnings and errors from each node.

The \texttt{PipelineStateValidator} enforces invariants between nodes: for
example, it verifies that every partition in the translation phase has a
non-null \texttt{risk\_level} and \texttt{rag\_strategy} field assigned by the
preceding routing node. Violations raise structured errors that are captured by
the orchestrator rather than propagating as unhandled exceptions.

\textbf{The composite node design pattern.} Each of the eight nodes in the
orchestrator is a \emph{composite agent}: a thin facade that delegates to one
or more specialist sub-agents. This design emerged from a refactoring in Week 13
that reduced an eleven-node pipeline to eight nodes without losing any
functionality. The advantage is twofold: the orchestrator sees a clean,
manageable eight-node graph, while the sub-agents remain independently testable
and replaceable. Table~\ref{tab:node-map} summarises the mapping.

[TABLE 4.4 --- Caption: ``Composite node design: eight orchestrator nodes and
their specialist sub-agents.'' --- Columns: Node Name, Orchestrator Label,
Sub-Agents Called, Layer.]

\subsection{Input Phase: File Processing and Streaming (Nodes 1--2)}
\label{subsec:nodes-input}

\textbf{Node 1 --- \texttt{file\_process}.} The \texttt{FileProcessor} node
is the only node whose failure is fatal to the pipeline. It orchestrates three
sub-agents:

\begin{itemize}
  \item \texttt{FileAnalysisAgent}: reads each uploaded SAS file, extracts
    top-level metadata (file size, encoding, SAS version markers), and detects
    \texttt{\%INCLUDE} references.
  \item \texttt{CrossFileDependencyResolver}: builds a dependency graph over
    all uploaded files using the detected \texttt{\%INCLUDE} links and shared
    libref conventions. Strongly connected components are identified via
    NetworkX's \texttt{strongly\_connected\_components} algorithm, flagging
    circular dependencies for manual review.
  \item \texttt{RegistryWriterAgent}: persists all file metadata and dependency
    edges to the SQLite file registry, enabling downstream agents to query
    cross-file context without re-reading the original files.
\end{itemize}

\textbf{Node 2 --- \texttt{streaming}.} The streaming node processes each file
through a producer--consumer pipeline: the \texttt{StreamAgent} reads the SAS
source character by character and pushes tokens into an \texttt{asyncio.Queue};
the \texttt{StateAgent} consumes from the queue, maintaining an FSM with 23
states that tracks the current position within the SAS programme (inside a DATA
step, inside a PROC step, inside a macro definition, inside a string literal,
etc.). This FSM-based approach enables processing of files up to 10,000 lines
without loading the full source into memory, satisfying the streaming performance
requirement (NF-02).

[FIGURE 4.2 --- Caption: ``Streaming FSM: producer--consumer pipeline between
StreamAgent and StateAgent via asyncio.Queue.'' --- Suggested: two boxes
(StreamAgent / StateAgent) connected by a Queue, with FSM state labels on the
StateAgent side.]

\subsection{Partitioning Phase: Boundary Detection and RAPTOR Clustering (Nodes 3--4)}
\label{subsec:nodes-partition}

\textbf{Node 3 --- \texttt{chunking}.} Boundary detection transforms the
token stream into a list of semantically coherent \texttt{PartitionIR} objects.
The \texttt{BoundaryDetectorAgent} implements a two-stage approach:

\begin{enumerate}
  \item \textbf{Deterministic detection} (approximately 80\% coverage): a lark
    grammar and a set of regex rules identify DATA-step and PROC-step
    boundaries deterministically. RUN statements, QUIT statements, and
    semicolons at the procedure top level are the primary boundary markers.
  \item \textbf{LLM fallback} (approximately 20\% of cases): when the
    deterministic rules cannot resolve an ambiguous boundary --- typically
    within macro-generated code or inside deeply nested \texttt{\%IF} blocks
    --- the \texttt{LLMBoundaryResolver} sends the ambiguous code excerpt to
    GPT-4o-mini with a structured boundary-detection prompt. The response is
    validated against the expected JSON schema before being accepted.
\end{enumerate}

The \texttt{PartitionBuilderAgent} converts the detected boundaries into
\texttt{PartitionIR} objects, each with a UUID \texttt{block\_id}, line range
(\texttt{line\_start}/\texttt{line\_end}), partition type
(one of ten \texttt{PartitionType} values), and an initially empty
\texttt{risk\_level} field.

\textbf{Node 4 --- \texttt{raptor}.} The RAPTOR node builds a hierarchical
retrieval tree over the partitioned blocks:

\begin{enumerate}
  \item \texttt{NomicEmbedder}: computes 768-dimensional embeddings for each
    partition using the Nomic Embed v1.5 model loaded via
    \texttt{sentence-transformers}. The model is loaded lazily on the first
    call (cold-start cost: approximately 2 seconds on CPU).
  \item \texttt{GMMClusterer}: fits a \ac{GMM} over the embeddings using
    scikit-learn's \texttt{GaussianMixture}. The number of components is
    set to $\min(n\_partitions / 2, 8)$ to avoid the degenerate case where
    $n\_components > n\_samples$, which would cause a rank-deficiency error.
  \item \texttt{ClusterSummarizer}: for each GMM cluster, sends the member
    partition snippets to the LLM with a summarisation prompt. The resulting
    cluster summary (e.g., ``BY-group running totals with RETAIN reset across
    three datasets'') is stored as a synthetic document in LanceDB alongside
    the leaf-level partitions.
  \item \texttt{RAPTORTreeBuilder}: assembles the leaf nodes and cluster
    summaries into a two-level tree (leaf $\rightarrow$ cluster; a root summary
    is added if the number of clusters $\geq 2$) and writes the RAPTOR back-links
    (\texttt{raptor\_leaf\_id}, \texttt{raptor\_cluster\_id},
    \texttt{raptor\_root\_id}) to each \texttt{PartitionIR}.
\end{enumerate}

\textbf{RAPTOR ablation results.} A controlled ablation study compared RAPTOR
hierarchical retrieval against a flat leaf-only index on the 50-file gold-standard
corpus. Table~\ref{tab:raptor-ablation-ch4} reports the results.

\begin{table}[h]
\centering
\caption{RAPTOR vs.\ flat-index retrieval on the gold-standard corpus.}
\label{tab:raptor-ablation-ch4}
\begin{tabular}{lcccc}
\hline
\textbf{Metric} & \textbf{Flat index} & \textbf{RAPTOR} & \textbf{Delta} & \textbf{Segment} \\
\hline
Hit-Rate@5 (all)     & 0.71 & 0.84 & +18\% & All \\
MRR (all)            & 0.54 & 0.63 & +17\% & All \\
Hit-Rate@5 (HIGH)    & 0.52 & 0.71 & +36\% & HIGH only \\
\hline
\end{tabular}
\end{table}

The RAPTOR advantage was most pronounced on HIGH-complexity partitions ($+36\%$),
where flat nearest-neighbour search fails to surface cluster-level context. On
MODERATE partitions the gap was $+12$ percentage points; LOW partitions showed
no meaningful difference, confirming that the cluster-level summaries add value
only when the partition requires multi-construct context.

\subsection{Routing, Persistence, Translation, and Merge (Nodes 5--8)}
\label{subsec:nodes-routing-merge}

\textbf{Node 5 --- \texttt{risk\_routing}.} The \texttt{RiskRouter} node
assigns a risk level and RAG strategy to each partition:

\begin{itemize}
  \item \texttt{ComplexityAgent}: extracts lexical and structural features from
    each partition (token count, nesting depth, macro references, BY clauses,
    hash object usage) and runs the Platt-scaled logistic regression classifier.
    The output is a probability vector over the four risk levels.
  \item \texttt{StrategyAgent}: maps the risk level to a RAG strategy:
    LOW $\rightarrow$ Static RAG; MODERATE + cross-file deps $\rightarrow$
    Graph RAG; HIGH / UNCERTAIN $\rightarrow$ Agentic RAG.
\end{itemize}

\textbf{Node 6 --- \texttt{persist\_index}.} Two sub-agents run concurrently:
the \texttt{PersistenceAgent} writes each \texttt{PartitionIR} to the SQLite
\texttt{conversion\_stages} table; the \texttt{IndexAgent} builds the NetworkX
dependency graph over all partitions, runs SCC analysis, and writes the
SCC membership to each partition's metadata.

\textbf{Node 7 --- \texttt{translation}.} The \texttt{TranslationPipeline}
is the most complex node. Its internal structure is detailed in
Chapter~\ref{ch:implementation}; at the architecture level, it exposes a single
\texttt{async process(state)} interface that receives the partitioned,
risk-routed state and returns the state augmented with Python translations,
verification results, and SemantiCheck Scores for each partition.

\textbf{Node 8 --- \texttt{merge}.} The \texttt{MergeAgent} assembles the
individual partition translations into a coherent Python script. The
\texttt{ImportConsolidator} deduplicates and sorts all \texttt{import} and
\texttt{from} statements from every partition's translated output. The
\texttt{DependencyInjector} orders the translated code blocks according to the
topological sort of the NetworkX dependency graph, ensuring that upstream
datasets are always computed before downstream references. The
\texttt{ReportAgent} generates the HTML/Markdown conversion report, including
the SemantiCheck Score breakdown, Z3 verification outcomes, and per-partition
risk levels.

% ---------------------------------------------------------------------------
\section{Data Architecture}
\label{sec:data-architecture}
% ---------------------------------------------------------------------------

\codara{} uses four distinct storage systems, each optimised for a different
access pattern. The rationale for this multi-store approach is separation of
concerns: a single database that tries to serve as an operational store, a
vector index, an analytical warehouse, and a distributed checkpoint store
simultaneously would be poorly optimised for all four.

[FIGURE 4.3 --- Caption: ``Multi-store data architecture: SQLite (operational),
LanceDB (vector), DuckDB (analytical), Redis (checkpoints), and Azure Blob/Queue
(cloud I/O).'' --- Suggested: architecture diagram with five storage boxes,
each labelled with their role and the pipeline components that access them.]

\subsection{SQLite API Database}
\label{subsec:sqlite-schema}

SQLite was chosen for the operational API database because it requires no server
process, is trivially deployable alongside the FastAPI application in Docker,
and is sufficient for the expected write load (one conversion at a time per
deployment in the initial release). SQLAlchemy 2.0 provides the ORM layer;
Write-Ahead Logging (WAL) mode is enabled at startup to reduce lock contention
under concurrent reads.

The schema comprises eight tables: \texttt{users}, \texttt{conversions},
\texttt{conversion\_stages} (one row per pipeline stage per conversion),
\texttt{kb\_entries}, \texttt{kb\_changelog}, \texttt{audit\_logs},
\texttt{corrections}, and \texttt{notifications}. Foreign key enforcement is
enabled via SQLite's \texttt{PRAGMA foreign\_keys = ON}.

[TABLE 4.5 --- Caption: ``SQLite schema overview: tables, key columns, and
access patterns.'' --- Columns: Table, Key Columns, Primary Access Pattern,
Row Estimate.]

\subsection{LanceDB Knowledge Base}
\label{subsec:lancedb-schema}

The knowledge base contains 330 verified SAS--Python translation pairs stored
as a LanceDB table named \texttt{sas\_python\_examples}. Each row is a 16-field
record combining the source text, target text, a 768-dimensional Nomic embedding
vector, and nine metadata fields: \texttt{complexity\_tier} (LOW/MOD/HIGH),
\texttt{failure\_mode} (one of six SAS error classes), \texttt{verified}
(boolean, true for gold-standard-derived pairs), \texttt{verification\_score}
(float, SemantiCheck Score at ingestion time), \texttt{version} (monotonically
increasing integer for rollback), \texttt{superseded\_by} (UUID pointer to the
newer version of this pair, if any), \texttt{category} (DATA step / PROC SQL /
RETAIN / macro / hash / etc.), \texttt{confidence} (LLM self-reported confidence
at generation time), and \texttt{source} (gold-standard / KB-expansion /
human-correction).

An IVF (Inverted File Index) index over the embedding column is built on
first query and cached on disk. Approximate nearest-neighbour search with cosine
similarity returns the top-$k$ candidates in under 5 ms for $k \leq 20$.

\subsection{DuckDB Audit Store}
\label{subsec:duckdb-schema}

DuckDB was chosen for the LLM audit log because its columnar storage and
vectorised query engine make analytical queries (e.g., ``average latency by
model in the last 24 hours'', ``token cost breakdown by risk level'') fast
without a separate warehouse infrastructure. The file \texttt{analytics.duckdb}
is opened in read-write mode by the \texttt{LLMAuditLogger} singleton during
pipeline execution.

Three tables are maintained: \texttt{conversion\_results} (one row per
translated partition, storing the python code, LLM confidence, failure mode
flag, model used, and KB examples retrieved), \texttt{llm\_calls} (one row
per LLM API call, with prompt hash, latency, estimated token cost, and success
flag), and \texttt{kb\_changelog} (KB version history, mirroring the SQLite
\texttt{kb\_changelog} table for analytical queries).

\subsection{Redis Checkpoint Store}
\label{subsec:redis-schema}

Redis serves as the LangGraph checkpoint backend. Every 50 processed partitions,
the \texttt{RedisCheckpointManager} serialises the full \texttt{PipelineState}
to a Redis key with a 24-hour TTL. If the pipeline crashes (OOM, network timeout,
LLM rate limit), the next invocation with the same \texttt{conversion\_id}
detects the existing checkpoint via \texttt{find\_checkpoint()} and resumes from
the last saved state, skipping all partitions already translated.

When Redis is unavailable, the pipeline continues in degraded mode without
checkpointing, logging a warning to the stage record. This fallback is
intentional: a Redis outage should not block translation, only disable
fault recovery.

\subsection{Azure Blob and Queue Storage}
\label{subsec:azure-storage-schema}

Two Azure storage services provide the cloud I/O layer in production deployments:

\textbf{Azure Blob Storage} (via \texttt{BlobStorageService}): SAS source files
uploaded through the API are stored as blobs under a container named
\texttt{sas-uploads/\{conversion\_id\}/}. The pipeline worker downloads the
blob to a temporary local path before processing. On completion, the merged
Python output and report are uploaded to
\texttt{sas-outputs/\{conversion\_id\}/}. When \texttt{AZURE\_STORAGE\_CONNECTION\_STRING}
is absent, the service falls back transparently to the local
\texttt{backend/uploads/} directory.

\textbf{Azure Queue Storage} (via \texttt{PipelineQueueService}): conversion
jobs are enqueued as JSON payloads (base64-encoded) with a
\texttt{visibility\_timeout} of 300 seconds. A daemon consumer thread polls
the queue every 2 seconds and dispatches jobs to \texttt{run\_pipeline\_sync}.
A poison-message guard dead-letters any message that has been dequeued five or
more times without successful processing. When the queue is not configured,
\texttt{FastAPI BackgroundTasks} is used as a single-process fallback.

% ---------------------------------------------------------------------------
\section{API and Service Layer Design}
\label{sec:api-design}
% ---------------------------------------------------------------------------

\subsection{Layered Architecture: Routes $\to$ Services $\to$ Core}
\label{subsec:layered-api}

The API follows a three-layer architecture that decouples HTTP concerns from
business logic and infrastructure:

\begin{enumerate}
  \item \textbf{Routes layer} (\texttt{backend/api/routes/}): six FastAPI
    routers handle HTTP request parsing, authentication, response serialisation,
    and error mapping. Routes call services; they contain no business logic.
  \item \textbf{Services layer} (\texttt{backend/api/} extracted in Week 15):
    three service modules provide infrastructure abstractions:
    \texttt{pipeline\_service.py} (\texttt{run\_pipeline\_sync}: downloads the
    SAS file, executes the 8-node pipeline, cleans up temporary files);
    \texttt{blob\_service.py} (\texttt{BlobStorageService}: file I/O abstraction
    with Azure/local fallback); \texttt{queue\_service.py}
    (\texttt{PipelineQueueService}: job dispatch abstraction with Azure/
    BackgroundTasks fallback).
  \item \textbf{Core layer} (\texttt{backend/partition/}): the LangGraph
    pipeline, all agents, and the data models. Completely independent of FastAPI;
    can be invoked directly via the CLI script
    \texttt{scripts/ops/run\_pipeline.py}.
\end{enumerate}

This separation was introduced in Week 15 after profiling revealed that the
routes were directly instantiating pipeline objects, making unit testing of
routing logic impossible without spinning up the full pipeline. The service
abstraction allows routes to be tested with mock services and the pipeline to
be tested independently.

\subsection{Key API Endpoints}
\label{subsec:key-endpoints}

[TABLE 4.6 --- Caption: ``Key REST API endpoints with method, path, authentication,
and description.'' --- Columns: Method, Path, Auth Required, Description.]

The system exposes 28 documented endpoints across six routers. The critical
path for a conversion involves four calls: \texttt{POST /api/conversions/upload}
(file upload, returns file UUIDs), \texttt{POST /api/conversions/start}
(triggers pipeline, returns conversion ID), repeated
\texttt{GET /api/conversions/\{id\}} (polls stage-level status at 1,200 ms
intervals), and \texttt{GET /api/conversions/\{id\}/download} (retrieves
the ZIP archive when complete).

\subsection{Authentication and Authorisation Design}
\label{subsec:auth-design}

All endpoints except \texttt{GET /api/health} require a valid JWT token in the
\texttt{Authorization: Bearer} header. Tokens are signed with HS256 using the
\texttt{CODARA\_JWT\_SECRET} environment variable, with a 24-hour expiry.
Passwords are hashed with bcrypt (12 rounds) via \texttt{passlib}.

The system supports three roles: \texttt{admin} (full access including user
management and audit logs), \texttt{user} (conversion and KB read/write), and
\texttt{viewer} (read-only access to conversions and analytics). Role checks
are implemented as FastAPI dependency injectors applied at the router level.

On first boot, if no users exist in the database, the system seeds two default
accounts using passwords generated via \texttt{secrets.token\_urlsafe(18)} and
printed to stdout. These can be pinned to specific values via the
\texttt{CODARA\_ADMIN\_PASSWORD} and \texttt{CODARA\_USER\_PASSWORD} environment
variables. The login and register endpoints are protected by a rate limiter
(5 attempts per IP per 60 seconds, implemented via an in-memory counter with
a 60-second rolling window) to mitigate credential-stuffing attacks.

% ---------------------------------------------------------------------------
\section{Gold Standard Corpus Construction}
\label{sec:corpus-construction}
% ---------------------------------------------------------------------------

\subsection{Motivation and Design Requirements}
\label{subsec:corpus-design}

A rigorous evaluation of any translation system requires a \emph{gold standard}:
a curated set of (source, expected-output) pairs whose correctness is
established independently of the system under evaluation. Without a gold
standard, evaluation degrades to ``the system produced output that looks
reasonable'', which is insufficient for scientific claims.

Our gold standard serves three distinct purposes: (1) training data for the
complexity classifier; (2) the retrieval corpus seed for the LanceDB knowledge
base (the top 45 pairs, verified by hand); and (3) the evaluation dataset for
boundary detection accuracy, RAPTOR retrieval quality, and end-to-end
translation success. The design requirement was therefore that the corpus cover
all four risk levels, all ten partition types, and all six known SAS failure
modes --- not just the easy cases.

\subsection{Construction Methodology}
\label{subsec:corpus-methodology}

The corpus was built in three tiers reflecting difficulty:

\textbf{Tier 1 --- Basic (15 pairs, prefix \texttt{gs\_}).} Each pair consists
of a self-contained SAS snippet (10--50 lines) illustrating a single construct
(DATA step, RETAIN, MERGE, FIRST./LAST., simple ETL, multi-output, hash object,
PROC MEANS, PROC FREQ, PROC SQL, macro, DO loop, \texttt{\%INCLUDE},
FILENAME statement), paired with a hand-written Python translation verified for
semantic correctness on three synthetic DataFrames.

\textbf{Tier 2 --- Medium (20 pairs, prefix \texttt{gsm\_}).} Longer SAS
programmes (50--150 lines) that combine multiple constructs: financial summary
with RETAIN and currency formatting, customer segmentation with PROC SQL and
macro parameters, clinical data cleaning with FIRST./LAST. and missing-value
logic, and so on. Translations were generated by GPT-4o and reviewed by the
project author.

\textbf{Tier 3 --- Hard (15 pairs, prefix \texttt{gsh\_}).} Enterprise-scale
SAS programmes (150--400 lines) with macro frameworks, cross-file dependencies,
and complex PROC SQL patterns. These pairs exercise the full agentic RAG path.
Translations required multiple LLM calls and manual correction passes.

Each \texttt{.gold.json} annotation file specifies the expected partition types,
boundary line numbers, complexity tier, and a list of the SAS failure modes
present in the source.

\subsection{Statistics and Limitations}
\label{subsec:corpus-statistics}

[TABLE 4.7 --- Caption: ``Gold standard corpus statistics by tier.'' ---
Columns: Tier, Pairs, Lines (total), Partitions (total), Complexity Distribution
(LOW/MOD/HIGH \%), Failure Modes Covered.]

The corpus totals 50 annotated SAS files (45 with hand-verified Python
translations), 721 individual blocks after boundary detection, and approximately
8,000 lines of SAS source. It is deliberately biased toward MODERATE and HIGH
complexity (60\% of pairs) because LOW-complexity translation is already reliably
handled by deterministic rules. The primary limitation is size: 45 pairs is a
small dataset by machine learning standards, and the complexity classifier was
therefore trained on the full 330-pair knowledge base rather than the
gold-standard corpus alone.

% ---------------------------------------------------------------------------
\section{Conclusion}
\label{sec:ch4-conclusion}
% ---------------------------------------------------------------------------

This chapter documented the architectural DNA of \codara{}: the empirical
process by which each major technology was selected, the composite node design
that organises eight specialised agents into a coherent LangGraph pipeline, the
multi-store data architecture that matches each storage pattern to its ideal
engine, the layered API design that decouples HTTP from business logic, and the
gold-standard corpus that provides the evaluative backbone for all claims in
Chapter~\ref{ch:evaluation}.

Two architectural decisions deserve particular emphasis because they were
counter-intuitive. First, the choice to use four distinct storage systems
(SQLite + LanceDB + DuckDB + Redis) rather than a single general-purpose
database was driven by measurement, not complexity preference: each store
outperforms the alternatives on its specific access pattern by a significant
margin. Second, the composite node pattern --- which adds an indirection layer
between the orchestrator and the specialist agents --- was essential for
testability: all 309 tests exercise individual sub-agents directly, without
requiring the full LangGraph pipeline to run.

The implementation details of every node and sub-agent are presented in the
following chapter.

% ---------------------------------------------------------------------------
% References for Chapter 4
% ---------------------------------------------------------------------------
%
% [1]  LangChain AI, "LangGraph: Build stateful multi-actor applications with
%      LLMs," https://langchain-ai.github.io/langgraph/, 2024.
% [2]  CrewAI Inc., "CrewAI: Framework for orchestrating role-playing,
%      autonomous AI agents," https://crewai.com, 2024.
% [3]  Wu, Q. et al., "AutoGen: Enabling Next-Gen LLM Applications via
%      Multi-Agent Conversation," in Proc. ICLR, 2024.
% [4]  Nygard, M. T., Release It! Design and Deploy Production-Ready Software,
%      2nd ed. Pragmatic Bookshelf, 2018.
% [5]  Nussbaum, Z. et al., "Nomic Embed: Training a Reproducible Long Context
%      Text Embedder," arXiv:2402.01613, 2024.
% [6]  LanceDB, "LanceDB: Developer-friendly, serverless vector database,"
%      https://lancedb.github.io/lancedb/, 2024.
% [7]  Chroma, "Chroma: The open-source embedding database,"
%      https://www.trychroma.com, 2024.
% [8]  Weaviate B.V., "Weaviate: Open-source vector database,"
%      https://weaviate.io, 2024.
% [9]  Johnson, J. et al., "Billion-scale similarity search with GPUs,"
%      IEEE Trans. Big Data, vol. 7, no. 3, pp. 535–547, 2021.
% [10] de Moura, L. and Bjørner, N., "Z3: An Efficient SMT Solver,"
%      in Proc. TACAS, LNCS 4963, 2008, pp. 337–340.
% [11] MacIver, D. and Hatfield-Dodds, Z., "Hypothesis: A new approach to
%      property-based testing," J. Open Source Softw., vol. 4, 2019.
